\block{How to deal with complex geometry in PINNs ?}{
    \warning \textbf{In practice :} Not so easy ! We need to find \textbf{\fcolorbox{color1!50}{color1}{how to sample in the geometry}}.

    \begin{minipage}{0.48\linewidth}
        \centering
        \begin{tcolorbox}[
            colback=color1!50, % Couleur de fond de la boîte
            colframe=color2, % Couleur du cadre de la boîte
            arc=2mm, % Rayon de l'arrondi des coins
            boxrule=2pt, % Épaisseur du cadre de la boîte
            breakable, enhanced jigsaw,
            width=\linewidth
            ]            
            \textbf{Solution :} Approach by levelset. \textcolor{red}{CITE}
        \end{tcolorbox}
    \end{minipage}	
    
    
    }
    
    \note[rotate=8, connection, width = 11cm, targetoffsetx=28cm, targetoffsety=5cm, roundedcorners=30, linewidth=1pt]{No mesh, so easy to go on complex geometry !}

    % \node[below left=2cm and 2cm] at (topright) {\includegraphics[width=2cm]{images/levelset/speaking.png}};
    
    % \begin{frame}{Limits}
    % 	\textbf{Claim on PINNs :} \textcolor{orange}{No mesh, so easy to go on complex geometry !}
	
% 	\warning \textit{In practice :} Not so easy ! We need to find \textcolor{orange}{how to sample in the geometry}.
	
% 	\vspace{8pt}
	
% 	\textbf{Solution :} Approach by levelset. \footnotesize\citep{sukumar_exact_2022}\normalsize
	
% 	\begin{center}
% 		\pgfimage[width=0.4\linewidth]{images/learn_levelset/levelset.png}
% 	\end{center}
	
% 	\vspace{5pt}

% 	\begin{center}
% 		\begin{minipage}{0.44\linewidth}
% 			\textbf{\textit{Advantages :}} \\
% 			\ding{217} Sample is easy in this case. \\
% 			\ding{217} Allow to impose in hard the BC :
% 			\vspace{-5pt}
% 			\begin{equation*}
% 				u_\theta(X)=\phi(X)w_\theta(X)+g(X)
% 			\end{equation*}
% 			(\ding{217} Can be used for $\phi$-FEM)
% 		\end{minipage}
% 		\begin{minipage}{0.44\linewidth}
% 			\textbf{\textit{Natural LevelSet :}} \\
% 			Signed Distance Function (SDF)
			
% 			\vspace{5pt}
% 			\textbf{\textit{Problem :}} SDF is a $\mathcal{C}^0$ function  \\
% 			$\Rightarrow$ its derivatives explode \\
% 			$\Rightarrow$ we \textcolor{orange}{need a regular levelset}
% 		\end{minipage}
% 	\end{center}
% \end{frame}

% \begin{frame}{Learn a regular levelset}	
% 	\vspace{-10pt}
% 	\begin{tcolorbox}[
% 		colback=other, % Couleur de fond de la boîte
% 		colframe=other, % Couleur du cadre de la boîte
% 		arc=2mm, % Rayon de l'arrondi des coins
% 		boxrule=0.5pt, % Épaisseur du cadre de la boîte
% 		breakable, enhanced jigsaw,
% 		width=\linewidth,
% 		opacityback=0.1
% 		]
		
% 		If we have a boundary domain $\Gamma$, the SDF is solution to the Eikonal equation:
		
% 		\begin{minipage}{0.7\linewidth}
% 			\hspace{100pt}
% 			$\left\{\begin{aligned}
% 				&||\nabla\phi(X)||=1, \; X\in\mathcal{O} \\
% 				&\phi(X)=0, \; X\in\Gamma \\
% 				&\nabla\phi(X)=n, \; X\in\Gamma
% 			\end{aligned}\right.$
% 		\end{minipage}
% 		\begin{minipage}{0.25\linewidth}
% 			\centering
% 			\pgfimage[width=0.7\linewidth]{images/learn_levelset/points_normals.png}
% 		\end{minipage}
		
% 		with $\mathcal{O}$ a box which contains $\Omega$ completely and $n$ the exterior normal to $\Gamma$.
% 	\end{tcolorbox}

% 	\textbf{How to do that ?} with a PINNs \footnotesize\citep{clemot_neural_2023}\normalsize ~by \textcolor{orange}{adding a regularization term},
% 	\vspace{-5pt}
% 	\begin{equation*}
% 		J_{reg} = \int_\mathcal{O} |\Delta\phi|^2
% 	\end{equation*}

% 	\begin{minipage}{0.32\linewidth}
% 		\centering
% 		\pgfimage[width=\linewidth]{images/learn_levelset/cat_levelset_loss.png}
% 	\end{minipage} 
% 	\begin{minipage}{0.32\linewidth}
% 		\centering
% 		\pgfimage[width=\linewidth]{images/learn_levelset/cat_levelset.png}
% 	\end{minipage} 
% 	\begin{minipage}{0.32\linewidth}
% 		\centering
% 		\pgfimage[width=\linewidth]{images/learn_levelset/cat_levelset_bc.png}
% 	\end{minipage} 

% 	\footnotesize
% 	\textit{Remark :} PINNs non parametric - \textcolor{red}{1 geometry}
% \end{frame}

% \begin{frame}{Poisson On Cat}	
% 	\ding{217} Solving (\ref{edp}) with $f=1$ (\textcolor{red}{non parametric}) and homogeneous Dirichlet BC ($g=0$). \\
% 	\ding{217} Looking for $u_\theta = \phi w_\theta$ with $\phi$ the levelset learned. 
	
% 	\begin{center}
% 		\begin{minipage}{0.32\linewidth}
% 			\textbf{Sampling}
			
% 			\vspace{-10pt}
% 			\centering
% 			\pgfimage[width=0.6\linewidth]{images/learn_levelset/cat_levelset_sampling.png}
% 		\end{minipage} \qquad 
% 		\begin{minipage}{0.32\linewidth}
% 			\centering
% 			\pgfimage[width=\linewidth]{images/learn_levelset/cat_poisson_loss.png}
% 		\end{minipage} 
% 	\end{center}

% 	\begin{center}
% 		\pgfimage[width=0.8\linewidth]{images/learn_levelset/cat_poisson_sol.jpg}
% 	\end{center}

% 	\footnotesize
% 	\textit{Remark :} Poisson on Bean \refappendix{frame:bean}
% \end{frame}